\documentclass[a4paper, serif, 10pt]{moderncv}

%% moderncv
% color and style
\moderncvcolor{black}
\moderncvstyle{banking}

%% page layout/margins
\usepackage[top=2.5cm, bottom=2.5cm, left=1.5cm, right=1.5cm]{geometry}

\nopagenumbers{}

%% Babel config for Norwegian language
% ref:
% - https://latex3.github.io/babel/guides/locale-norwegian.html
\usepackage[norsk]{babel}

%% fontspec
\usepackage{fontspec}

\IfFontExistsTF{Linux Libertine}{
  \setmainfont[Ligatures={TeX, Common}, Numbers=OldStyle]{Linux Libertine}
}{}
\IfFontExistsTF{Linux Libertine O}{
  \setmainfont[Ligatures={TeX, Common}, Numbers=OldStyle]{Linux Libertine O}
}{}

%% CTeX
% CJK support, used to show CJK characters in the resume
%
% - heading=false: disable ctex heading
% - punct=kaiming: use kaiming punctuations styles for CJK
% - scheme=plain: use plain scheme, do not override \normalsize font size
% - space=auto: space settings for CJK characters
%
% CJK fonts are resolved per locale language, e.g. Noto Serif/Sans CJK JP
% for Japanese resumes, with fallback to the next region if not installed
%
% ref:
% - http://ctan.mirrorcatalogs.com/language/chinese/ctex/ctex.pdf
\usepackage[UTF8, heading=false, punct=kaiming, scheme=plain, space=auto]{ctex}

\IfFontExistsTF{Noto Serif CJK SC}{
  \setCJKmainfont{Noto Serif CJK SC}
}{}
\IfFontExistsTF{Noto Sans CJK SC}{
  \setCJKsansfont{Noto Sans CJK SC}
}{}

%% line spacing
\usepackage{setspace}
\setstretch{1.125}

%% URL styling - use normal text instead of monospace
\urlstyle{same}

%% Patch \httplink and \httpslink to handle full URLs
% The original moderncv macros always prepend http:// or https:// to URLs.
% This causes issues when the URL already has a protocol (e.g., https://example.com)
% resulting in malformed URLs like https://https://example.com.
% This patch checks if the URL already contains "://" and uses it directly if so.
% Note: We use \str_set:Nx (x-expansion) to expand macros like \@homepage first.
\ExplSyntaxOn
\str_new:N \g_yamlresume_url_str

\RenewDocumentCommand{\httpslink}{O{}m}{
  \str_set:Nx \g_yamlresume_url_str {#2}
  \str_if_in:NnTF \g_yamlresume_url_str {://}
    { \url{#2} }
    { \url{https://#2} }
}

\RenewDocumentCommand{\httplink}{O{}m}{
  \str_set:Nx \g_yamlresume_url_str {#2}
  \str_if_in:NnTF \g_yamlresume_url_str {://}
    { \url{#2} }
    { \url{http://#2} }
}
\ExplSyntaxOff

\name{Ingrid Johansen}{}
\title{HR-sjef}
\phone[mobile]{+47 123 45 678}
\email{ingrid.johansen@example.no}
\homepage{https://www.linkedin.com/in/ingridjohansen}

\address{Karl Johans gate 15, Oslo, Oslo, Norge, 0154}{}{}

\extrainfo{{\small \faLinkedin}\ \href{https://www.linkedin.com/in/ingridjohansen}{@ingridjohansen} {} {} {} • {} {} {} 
{\small \faTwitter}\ \href{https://twitter.com/ingridhr}{@ingridhr}}

\begin{document}

\maketitle

\section{Grunnleggende informasjon}

\cvline{}{\begin{itemize}
\item Erfaren HR-leder med over 10 års erfaring fra teknologi- og finanssektoren
\item Sterk kompetanse innen rekruttering, organisasjonsutvikling og arbeidsrett
\item Dokumentert evne til å bygge sterke team og fremme en positiv arbeidskultur
\item Strategisk tenker med fokus på å koble HR-initiativer til forretningsmål
\end{itemize}}
\section{Utdanning}

\cventry{aug. 2012–juni 2014}
        {Master, Organisasjonspsykologi og ledelse, Poeng: 3.8}
        {Universitetet i Oslo}
        {\url{https://www.uio.no}}
        {}
        {\begin{itemize}
\item Utviklet dyp forståelse for sammenhengen mellom HR-praksis og organisasjonens resultater
\item Gjennomførte flere prosjekter knyttet til medarbeiderundersøkelser og kulturutvikling
\item Fikk solid innføring i norsk arbeidsrett og tariffsystemer
\end{itemize}
\textbf{Kurs}: Arbeids- og organisasjonspsykologi, Ledelse og organisasjonsutvikling, HR-strategi, Arbeidsrett, Rekruttering og seleksjon, Lønns- og belønningssystemer}
\section{Arbeidserfaring}

\cventry{jan. 2020–Nå}
        {HR-sjef}
        {Telenor}
        {\url{https://www.telenor.no}}
        {}
        {\begin{itemize}
\item Ansvarlig for HR-strategi og -drift for 2500 ansatte i Norge
\item Ledet arbeidet med å implementere ny HR-plattform (SAP SuccessFactors) og digitalisere HR-prosesser
\item Utviklet og implementerte talentprogrammer som økte intern rekruttering med 25\%
\item Rådgiver for ledergruppen innen organisasjonsutvikling, endringsledelse og arbeidsrett
\item Forhandlet frem nye tariffavtaler og sikret god dialogs med fagforeningene
\end{itemize}
\textbf{Nøkkelord}: HR-strategi, Talentutvikling, Endringsledelse, Arbeidsrett, Digitalisering}

\cventry{aug. 2016–des. 2019}
        {HR-rådgiver}
        {DNB}
        {\url{https://www.dnb.no}}
        {}
        {\begin{itemize}
\item Rekrutterte og onboardet over 100 kandidater årlig til roller innen IT og finans
\item Utviklet kompetanseutviklingsprogrammer og gjennomførte lederopplæring i samarbeid med eksterne leverandører
\item Veiledet ledere i personaloppfølging, konflikthåndtering og lønnsforhandlinger
\item Bidro til å redusere sykefraværet gjennom målrettet HMS-arbeid og oppfølging
\end{itemize}
\textbf{Nøkkelord}: Rekruttering, Kompetanseutvikling, Lønnsforhandlinger, HMS, Lederveiledning}
\section{Språk}

\cvline{Norsk}{Morsmål eller tospråklig nivå \hfill \textbf{Nøkkelord}: Morsmål}
\cvline{Engelsk}{Fullt profesjonelt nivå \hfill \textbf{Nøkkelord}: TOEFL 110, Forretningsengelsk}
\section{Ferdigheter}

\cvline{Rekruttering og seleksjon}{Ekspert \hfill \textbf{Nøkkelord}: Rekrutteringsstrategi, Intervjuteknikk, Assessment, Employer branding}
\cvline{Arbeidsrett og compliance}{Avansert \hfill \textbf{Nøkkelord}: Arbeidsmiljøloven, Tariffavtaler, GDPR, Arbeidskontrakter}
\cvline{Organisasjonsutvikling}{Avansert \hfill \textbf{Nøkkelord}: Endringsledelse, Medarbeiderundersøkelser, Kulturutvikling, Teamutvikling}
\cvline{HR-systemer og analyse}{Viderekommen \hfill \textbf{Nøkkelord}: SAP SuccessFactors, Workday, Excel, HR-analyse}
\section{Utmerkelser}

\cventry{nov. 2022}
        {Norsk HR-forening}
        {Årets HR-leder}
        {}
        {}
        {\begin{itemize}
\item Tildelt for fremragende arbeid med å implementere en ny HR-strategi i Telenor
\item Anerkjent for å ha økt medarbeidertilfredsheten med 20\% over to år
\end{itemize}}
\section{Sertifikater}

\cventry{mars 2021}
        {HR Certification Institute}
        {Sertifisert HR-profesjonell (SPHR)}
        {\url{https://www.hrci.org}}
        {}
        {}

\cventry{juni 2018}
        {International Project Management Association}
        {IPMA Level C - Sertifisert prosjektleder}
        {\url{https://www.ipma.world}}
        {}
        {}
\section{Publikasjoner}

\cventry{apr. 2023}
        {Fremtidens HR: Strategisk partnerskap i praksis}
        {HR-Nytt}
        {\url{https://www.hrnytt.no/artikkel/fremtidens-hr}}
        {}
        {\begin{itemize}
\item Artikkel om hvordan HR kan bli en strategisk partner i norske virksomheter
\item Diskuterer hvordan teknologi og data kan forbedre HR-prosesser
\end{itemize}}
\section{Prosjekter}

\cventry{jan. 2021–des. 2021}
        {Innføring av SAP SuccessFactors for 5000 ansatte}
        {Implementering av ny HR-plattform}
        {\url{https://www.telenor.no/hr-plattform}}
        {}
        {\begin{itemize}
\item Ledet prosjektet med å implementere ny HR-plattform i Telenor
\item Ansvarlig for behovsanalyse, leverandørvalg og utrulling
\item Oppnådde 30\% reduksjon i administrasjonstid for HR-prosesser
\end{itemize}
\textbf{Nøkkelord}: SAP SuccessFactors, Prosjektledelse, Endringsledelse}
\section{Interesser}

\cvline{Friluftsliv}{Fjellturer, Ski, Fisking}
\cvline{Litteratur}{Skjønnlitteratur, Biografier}
\section{Frivillig arbeid}

\cventry{sep. 2018–Nå}
        {Mentor}
        {Ungt Entreprenørskap}
        {\url{https://www.ue.no}}
        {}
        {\begin{itemize}
\item Fungerer som mentor for ungdomsskoleelever i entreprenørskapsprosjekter
\item Veileder team i forretningsidéutvikling og presentasjonsteknikk
\item Bidrar til å styrke unges kompetanse innen innovasjon og samarbeid
\end{itemize}}
    
\end{document}